\documentclass[10pt,a4paper]{amsart}
\usepackage[margin=19mm]{geometry}
\usepackage{amsmath,amssymb,mathtools}
\usepackage[scr=rsfs,cal=euler]{mathalfa}
\usepackage{xfrac}
\usepackage[unicode,hidelinks,bookmarks=false]{hyperref}
\newcommand{\ie}{i.e.\ }
\newcommand{\cf}{cf.\ }
\newcommand{\resp}{respectively\ }
\newcommand{\Subsection}{\subsection}
\providecommand{\nobreakdash}{\nobreak\hbox{-}}
\setlength{\emergencystretch}{2em}
\multlinegap0pt
\usepackage{amssymb}
\usepackage{booktabs}
\usepackage{xcolor}
\usepackage{tikz}
\usepackage{siunitx}
\newtheorem{assumption}{Assumption}
\newtheorem{lemma}{Lemma}
\title{\LARGE \bf Finite Memory Belief Approximation for Optimal Control in \\ Partially Observable Markov Decision Processes}
\author{Mintae Kim}
\date{Source: arXiv:2601.03132v1 (CC BY 4.0)}
\begin{document}
\maketitle

\noindent\emph{Source and licence.} \LARGE \bf Finite Memory Belief Approximation for Optimal Control in \\ Partially Observable Markov Decision Processes. Version arXiv:2601.03132v1 (CC BY 4.0), distributed by arXiv under the Creative Commons Attribution 4.0 International licence, http://creativecommons.org/licenses/by/4.0/, as recorded in the arXiv licence metadata for this exact version and shown on the abstract page https://arxiv.org/abs/2601.03132v1. Full licence notice: \texttt{licenses/arxiv-2601.03132-CC-BY-4.0.txt}.\par\smallskip

\noindent\textit{Editorial scope.} Counted unit: the article's lemma lem:cost-sensitivity (source0.tex lines 362--373). The article's complete original proof of it is delivered with it (source0.tex lines 375--477, one \texttt{proof} environment of 103 lines). No mathematics is written, repaired or re-derived: the statement and the proof are the article's own lines, in the article's own notation, and no retained line is substituted or re-flowed.  Retained uncounted as the source-local context the proof needs: lines 340--348; lines 350--357.  Nothing was omitted from any retained span, so the package carries no omission record.  No retained line contains a citation, so the wrapper carries no bibliography and no bibliography entry.  No retained line contains a figure environment, an image-inclusion command or drawing code, so no image is delivered and the package declares no asset dependency.  Editorial formatting only, in addition: an AMS \texttt{amsart} preamble that restates the article's own declarations and macros used by the retained lines, namely the environments \texttt{assumption}, \texttt{lemma} on separate counters, together with the standard packages those lines need.  The retained environments print in this document as Assumption 1, Lemma 1 in order of appearance, whereas the article prints the same items with the numbering of its own full document.  The complete native source of the article as distributed in the arXiv e-print for this exact version is bundled unchanged as \texttt{sources/192208/source0.tex}.
\begin{assumption}[Policy-conditional state regularity]
\label{assump:state-moment-regularity}
For a fixed admissible policy $\pi$, there exists a constant $M_\pi\in (0, \infty)$ for both measure $b_t^\pi$ and $\hat{b}_t^{(H),\pi}$ such that
\begin{equation}
\sup_{t\ge 0}\,
\mathbb{E}_{x\sim b_t^\pi}\!\left[\|x\|^2\right] \le M_\pi, \quad
\sup_{t\ge 0}\, \mathbb{E}_{x\sim \hat b_t^{(H),\pi}}\!\left[\|x\|^2\right] \le M_\pi.
\end{equation}
\end{assumption}

\begin{assumption}[Quadratic growth and smoothness of cost]
\label{assump:cost-regularity}
There exists a constant $K_c>0$ and $K_g>0$ such that for all $(x,u)\in\mathcal{X}\times\mathcal{U}$,
\begin{align}
|c(x,u)| \le K_c\left(1+\|x\|^2+\|u\|^2\right), \\
\|\nabla_x c(x,u)\| \le K_g\left(1+\|x\|+\|u\|\right).
\end{align}
\end{assumption}

\begin{lemma}[Belief-level cost sensitivity under $W_2$]
\label{lem:cost-sensitivity}
Suppose Assumptions~\ref{assump:state-moment-regularity} and~\ref{assump:cost-regularity} hold for a fixed policy $\pi$.
Then there exists a finite constant $L_\pi\in (0, \infty)$ such that for any $u\in\mathcal{U}$ and any $b,\tilde b\in\mathcal{P}_2(\mathcal{X})$ satisfying
$\mathbb{E}_{x\sim b}[\|x\|^2] \le M_\pi$ and
$\mathbb{E}_{x\sim \tilde b}[\|x\|^2] \le M_\pi$,
\begin{equation}
\left|\bar c(b,u)-\bar c(\tilde b,u)\right|
\le
L_\pi\left(1+\|u\|^2\right)W_2(b,\tilde b).
\end{equation}
\end{lemma}

\begin{proof}
Fix $u\in\mathcal{U}$ and $b,\tilde b\in\mathcal{P}_2(\mathcal{X})$ satisfying the stated second-moment bounds.
Let $(X,\tilde X)$ be any coupling of $(b,\tilde b)$.
Define $X_\lambda:=\tilde X+\lambda(X-\tilde X)$ for $\lambda\in[0,1]$.
By the fundamental theorem of calculus,
\begin{equation}
c(X,u)-c(\tilde X,u)
=
\int_0^1 \nabla_x c(X_\lambda,u)^\top (X-\tilde X)\,d\lambda.
\end{equation}
Taking absolute values and applying Cauchy-Schwarz yields
\begin{equation}
|c(X,u)-c(\tilde X,u)|
\le
\left(\int_0^1 \|\nabla_x c(X_\lambda,u)\|\,d\lambda\right)\|X-\tilde X\|.
\end{equation}
Taking expectation and applying Cauchy-Schwarz gives
\begin{multline}
\mathbb{E}\!\left[|c(X,u)-c(\tilde X,u)|\right]
\le \\
\left(
\mathbb{E}\!\left[\left(\int_0^1 \|\nabla_x c(X_\lambda,u)\|\,d\lambda\right)^2\right]
\right)^{\!1/2}
\left(\mathbb{E}\|X-\tilde X\|^2\right)^{\!1/2}.
\end{multline}
By Assumption~\ref{assump:cost-regularity},
\begin{equation}
\|\nabla_x c(X_\lambda,u)\|
\le
K_g\left(1+\|X_\lambda\|+\|u\|\right).
\end{equation}
Moreover, $\|X_\lambda\|\le \|\tilde X\|+\|X-\tilde X\|$ implies
\begin{equation}
1+\|X_\lambda\|+\|u\|
\le
1+\|\tilde X\|+\|X-\tilde X\|+\|u\|.
\end{equation}
By Jensen's inequality,
\begin{equation}
\left(\int_0^1 \|\nabla_x c(X_\lambda,u)\|\,d\lambda\right)^2
\le
\int_0^1 \|\nabla_x c(X_\lambda,u)\|^2\,d\lambda,
\end{equation}
and hence
\begin{multline}
\mathbb{E}\!\left[\left(\int_0^1 \|\nabla_x c(X_\lambda,u)\|\,d\lambda\right)^2\right]
\le \\
K_g^2\,
\mathbb{E}\!\left[
\int_0^1 \left(1+\|X_\lambda\|+\|u\|\right)^2 d\lambda
\right].
\end{multline}
Using $(a+b+c)^2\le 3(a^2+b^2+c^2)$ and the bound on $\|X_\lambda\|$,
\begin{multline}
\left(1+\|X_\lambda\|+\|u\|\right)^2
\le
3\left(1+\|u\|^2+\|X_\lambda\|^2\right)
\le \\
3\left(1+\|u\|^2+2\|\tilde X\|^2+2\|X-\tilde X\|^2\right).
\end{multline}
Therefore,
\begin{multline}
\label{eq:sec-4-eq-27}
\mathbb{E}\!\left[\left(\int_0^1 \|\nabla_x c(X_\lambda,u)\|\,d\lambda\right)^2\right]
\le
3K_g^2\left(1+\|u\|^2\right)
+ \\
6K_g^2\,\mathbb{E}[\|\tilde X\|^2]
+
6K_g^2\,\mathbb{E}[\|X-\tilde X\|^2].
\end{multline}
Since $X\sim b$ and $\tilde X\sim \tilde b$ satisfy $\mathbb{E}[\|X\|^2] \le M_\pi$ and $\mathbb{E}[\|\tilde X\|^2] \le M_\pi$, we also have
\begin{equation}
\label{eq:sec-4-eq-28}
\mathbb{E}[\|X-\tilde X\|^2] \le 2\mathbb{E}[\|X\|^2] + 2\mathbb{E}[\|\tilde X\|]^2 \le 4M_\pi.
\end{equation}
Combining (\ref{eq:sec-4-eq-27}) and (\ref{eq:sec-4-eq-28}) yields the uniform bound
\begin{equation}
\mathbb{E}\!\left[\left(\int_0^1 \|\nabla_x c(X_\lambda,u)\|\,d\lambda\right)^2\right]
\le
\tilde L_\pi^2\left(1+\|u\|^2\right),
\end{equation}
where one may take $\tilde L_\pi:=K_g\sqrt{3+24M_\pi}$.
Substituting into the Cauchy-Schwarz bound gives
\begin{multline}
\mathbb{E}\!\left[|c(X,u)-c(\tilde X,u)|\right]
\le \\
\tilde L_\pi\left(1+\|u\|^2\right)^{1/2}\left(\mathbb{E}[\|X-\tilde X\|^2]\right)^{1/2}.
\end{multline}
Since $\left|\bar c(b,u)-\bar c(\tilde b,u)\right|\le \mathbb{E}\!\left[|c(X,u)-c(\tilde X,u)|\right]$ and $(1+\|u\|^2)^{1/2}\le 1+\|u\|^2$, we obtain
\begin{equation}
\left|\bar c(b,u)-\bar c(\tilde b,u)\right|
\le
\tilde L_\pi\left(1+\|u\|^2\right)\left(\mathbb{E}\|X-\tilde X\|^2\right)^{1/2}.
\end{equation}
Taking the infimum over all couplings $(X,\tilde X)$ yields
\begin{equation}
\left|\bar c(b,u)-\bar c(\tilde b,u)\right|
\le
\tilde L_\pi\left(1+\|u\|^2\right)W_2(b,\tilde b).
\end{equation}
Setting $L_\pi:=\tilde L_\pi$ completes the proof.
\end{proof}

\end{document}
